\section{Introduction}\label{sec:introduction}
Network models often combine routing decisions with a choice among services
or operating modes. A route-specific adjustment to a service cost gives a
product of an arc flow and a service weight. Only a small fraction of the
possible route--service pairs may occur. Convexifying each product separately
discards restrictions imposed jointly by flow conservation and the common
simplex. The exact joint hull is available through classical simplex-vertex
disaggregation, but its usual formulation copies every arc for every state.
The question studied here is how the network and the retained products can
reduce that expansion, and when the resulting hull admits simple inequalities
in the original coordinates.

We consider a bounded equality-flow polytope
$P=\{x:Ax=b,\ 0\le x\le u\}$ on a directed multigraph and a vector
$y\ge0$ with $\sum_jy_j\le1$. Only products $z_{ej}=x_e y_j$ in a
specified set $\Obs$ are retained. Let $H$ be their joint convex hull.
Equality conservation is essential to the circulation arguments. Converting
balance inequalities to equalities with slack arcs is possible, but the graph
assumptions must then be checked on the transformed network. The general
network and block results permit arbitrary arc orientations, parallel arcs,
loops, zero simplex weights, and capacity degeneracy. The flat-chain results
use the specified directed unit-data topology, with arbitrary observations
and zero weights allowed.

The main reductions use two kinds of information. Circulations separate across
the cyclic blocks of the underlying undirected graph. Within each block,
states with no observed product can be merged exactly. More strongly, the
remaining freedom of an observed state consists precisely of circulations
supported on its unobserved arcs. Thus a graph with many cycles can still
have no missing state coordinates when its observations meet every cycle.
Conversely, a long chain can retain a small shared state profile even though
its cycle rank grows with its length.

Our contributions are organized around these two reductions and their limits.
\begin{enumerate}[leftmargin=*,itemsep=5pt]
\item \emph{Observation-sensitive formulations on arbitrary graphs.}
For a cyclic block $B$ and a label $j$ observed there, let $\rho_{Bj}$
be the cycle rank of the subgraph formed by arcs without an observation in
state $j$, including its isolated vertices when counting components.
The construction in \cref{thm:observed-compression} uses exactly
$\sum_{B,j}\rho_{Bj}$ additional coordinates, with unit flow, product,
and auxiliary coefficients in a specified rational row scaling.
This is the residual-coordinate count of the construction; it is not a
minimum over all extended formulations. Forest complements give a direct
original-coordinate hull. The same count is the minimum number of additional
\emph{individual products} needed to reconstruct each state on the ambient
circulation space, a narrower completion statement.

\item \emph{Explicit separation and recovery from graph structure.}
Cycles and theta blocks admit interval and six-support descriptions with
unit flow/product coefficients. Any number of parallel paths admits a
transportation cut description and a flow separator. For arbitrary blocks
of maximum cycle rank $r$, finite support libraries give exact separation
in $2^{O(r^2)}N$ rational arithmetic operations and compact recovery in
$2^{O(r^3)}N$, where $N=|V|+|E|+m+|\Obs|$, after observation grouping.
Neither procedure calls an online LP. The coefficient bound is
$H_r$ from \eqref{eq:rank-coefficient-bound}; in particular $H_3=2$ is
sharp on a unit-capacity $K_4$ with two explicit labels and five products.
The recovery parameter bound is deliberately distinguished from the smaller
separation bound.

\item \emph{A sharp observed-label threshold on flat chains.}
On a chain of parallel arc pairs with one bypass, the full shared profile
can be eliminated using a finite library depending only on the number $a$
of labels observed anywhere on the chain. After local checks and bound
grouping, at most two observed labels need five profile tests and explicit
interval recovery. Three labels use
16 positive circuits; two exceptional bypass coefficients can be repaired
using one flow-balance equation. Consequently every observation pattern
with $a\le3$ has a unit flow/product description, whereas a five-vertex,
nine-arc example with four observed labels does not. This threshold is
independent of the number of unused simplex coordinates and of chain length.

\item \emph{Coefficient obstructions in sparse original coordinates.}
A Fibonacci family forces exponentially growing product-coefficient ratios
on series--parallel networks with only linearly many observed products;
the obstruction persists on simple graphs of maximum degree three.
For unrestricted networks, a coordinate-section transfer of classical
three-way transportation universality gives arbitrary coefficient ratios
already with two explicit labels and unit network data. Its large
coefficient magnitudes coexist with polynomial encoding and polynomial-size
extended formulations. The coordinate-section proofs retain the relevant
products as actual coordinates of the sparse hull, rather than relying on
projection from a larger product space.
\end{enumerate}

\begin{table}[tb]\centering\small
\caption{Structural scope of the positive results. A unit coefficient is in
$\{0,\pm1\}$ on flow and product coordinates; simplex coefficients may contain
the rational network data. Bounds use the stated row scaling, before any
clearing of simplex denominators.}\label{tab:structural-scope}
\begin{tabular}{>{\raggedright\arraybackslash}p{.25\linewidth}
                >{\raggedright\arraybackslash}p{.44\linewidth}
                >{\raggedright\arraybackslash}p{.20\linewidth}}\toprule
Structure & Exact representation or procedure & Main locator\\\midrule
Arbitrary equality-flow graph & Observation-sensitive EF with
$\sum_{B,j}\rho_{Bj}$ additional coordinates & \Cref{thm:observed-compression}\\[3pt]
Unobserved arcs form forests in each observed block/state & Unit
original-coordinate formulation & \Cref{cor:forest-hull}\\[3pt]
Cycle or theta blocks & Unit interval/support cuts; compact recovery &
\Cref{thm:theta-hull,prop:theta-recovery}\\[3pt]
Parallel-path blocks & Complete subset cuts; transportation-flow separator &
\Cref{thm:parallel-hull}\\[3pt]
Block cycle rank at most $r$ & Coefficients bounded by $H_r$;
finite support and recovery libraries & \Cref{thm:bounded-rank,thm:rank-recovery}\\[3pt]
Flat chain, $a$ observed labels & Fixed-$a$ profile oracle; unit coefficients
for $a\le3$, sharp at four & \Cref{cor:chain-observed-labels}\\\bottomrule
\end{tabular}
\end{table}

These results refine established convexification methods. Simplex
disaggregation, shared-simplex gluing, reduced reformulation--linearization,
transportation projections, and normal-fan refinements supply the underlying
mechanisms. \Cref{subsec:prior-work} gives precise comparisons, especially
with the network convexification of \citet{KhademniaDavarnia2025}.
The claims concern the particular sparse formulations, complete structural
oracles, and coefficient boundaries established here.

Exact rational implementations and numerical LP formulations make the
distinction between these outputs concrete. The experiments compare with
full disaggregation, global removal of unused labels, native fixed-weight
state-flow LPs, and elementary one-flow or independent-network reductions
where applicable. The strongest globally merged LP is often fastest.
When every label is observed somewhere but each block sees few labels,
local compression still reduces 7,215 state-flow variables to 495 total
variables, or 367 after observed elimination. In that synthetic control,
initial compression reduces the median total time from about 30 to 10
milliseconds; elimination gives the smallest model but takes longer.
The dedicated exact oracle also loses to a simple LP on some long-chain
queries. These results support model-size and certificate benefits without
a universal runtime claim.

\Cref{sec:foundations,sec:blocks} establish the model and known disaggregation
and gluing facts. \Cref{sec:compression,sec:structured-oracles,sec:bounded-rank}
develop observation-sensitive formulations and graph-structured oracles.
\Cref{sec:universality,sec:sp-growth} prove the coefficient obstructions,
and \cref{sec:fixed-chains} gives the complementary fixed-label chain results.
\Cref{sec:computation} reports implementations, verification, experiments,
and a two-supplier service-transportation interpretation.
